\section*{Project contract}
Use the reflected 32-bit CRC specified below. The two update entries take a full byte
pointer in \texttt{a0}, a length \(0\leq n\leq4096\) in \texttt{a1}, and the
current raw unsigned state in \texttt{a2}. Its upper 32 bits are zero. Return the next
raw state, also zero-extended, in \texttt{v0}. The bit entry is
\texttt{crc32\_bit\_update}; \texttt{crc32\_table\_update} also receives a
word-aligned pointer to 256 prepared words in \texttt{a3}. Source bytes may be unaligned.
Read only the selected source interval and, for the table entry, the supplied 1024-byte
table. Write no caller bytes. Zero input length returns the state unchanged and reads
neither source nor table. The table and source slices are disjoint.

Implement \texttt{crc32\_table\_init} separately. It takes a word-aligned table
pointer in \texttt{a0}, initializes exactly 256 words and returns 256 in
\texttt{v0}. Its result must not depend on the old table contents. Only those
1024 caller bytes may be written. Submission-owned constants and private stack storage
are allowed in all routines. Functions receive full pointers, including addresses above
32 bits; adding a small offset must retain the full base.

The file entry, \texttt{crc32\_file}, receives a terminated path in \texttt{a0},
a scratch pointer in \texttt{a1}, a capacity \(1\leq c\leq256\) in
\texttt{a2}, and method zero or one in \texttt{a3}. Zero selects the bit update;
one selects the table update. Return the finalized unsigned checksum in \texttt{v0}
and status zero in \texttt{v1}. An unsuccessful open returns zero with status one;
a read error returns zero with status two. A successfully opened descriptor must be
closed on both success and read failure. An unsuccessful open creates no descriptor
to close. Do not return a partial checksum as a successful file result.

The file routine owns the selected scratch capacity and the first 1024 bytes of its
\texttt{crc\_table} workspace. It may read and write within them, but only bytes
reported by successful reads may contribute to the checksum. Preserve the path and
all other caller bytes. Retain the supplied arenas and caller labels for grading;
independent kernel calls may select offsets within the larger backing arenas.

Every entry preserves all bits of \texttt{s0} through \texttt{s7}, \texttt{fp}
and \texttt{sp}, and returns through the supplied \texttt{ra}. The incoming stack
is eight-byte aligned. The kernels perform no I/O. The file entry opens the requested
path read-only and produces no console output. Each file call begins a new checksum;
only the explicit raw-state argument carries state between kernel calls.
